\section{Experiencia}
\resumeSubHeadingListStart
\resumeExp{Ingeniero de Software Junior}{Reispsolutions S.R.L.}{Jul 2024 -- Jul 2026}{Santo Domingo, Rep. Dom.}
\resumeItemListStart
\resumeItem{Diseñé la arquitectura de un Sistema de Gestión de Aprendizaje (LMS) distribuido en ASP.NET junto a ingenieros senior, estandarizando la capacitación del personal en decenas de restaurantes multimarca.}
\resumeItem{Desarrollé servicios backend en producción con .NET 9 aplicando Command Query Responsibility Segregation (CQRS), Domain-Driven Design (DDD) y Entity Framework Core, realizando revisiones de código entre pares.}
\resumeItem{Desarrollé una plataforma web integral de gestión clínica que automatizó el registro de pacientes, la agenda de citas y recordatorios por SMS/correo, reduciendo drásticamente la coordinación manual en recepción.}
\resumeItem{Colaboré en equipos multifuncionales para crear aplicaciones de quiosco de autoservicio en tabletas Android con Jetpack Compose, optimizando el flujo de pedidos en múltiples franquicias de restaurantes.}
\resumeItemListEnd

\resumeExp{Asistente de Manufactura y Telemetría}{Equipo NASA Human Exploration Rover Challenge (HERC) de INTEC}{Ago 2024 -- Abr 2025}{Santo Domingo, Rep. Dom.}
\resumeItemListStart
\resumeItem{Colaboré con los subequipos mecánico y eléctrico en el diseño y fabricación de la carga útil de detección de líquidos del rover 2025: un actuador de husillo con motor paso a paso para posicionar sensores de agua.}
\resumeItem{Programé el sistema embebido del rover en ESP32-S3, decodificando señales PWM de un receptor de radiocontrol (RC) en una máquina de estados en tiempo real que controlaba la tracción y la ejecución de tareas.}
\resumeItem{Dirigí el desarrollo full-stack de la plataforma web oficial del equipo Apolo 27 y coordiné un equipo estudiantil de desarrollo para crear videojuegos interactivos de divulgación STEM en República Dominicana.}
\resumeItemListEnd
\resumeSubHeadingListEnd
